% New introductory reinforcement-learning exercise (TASKS.md D3); faculty draft.
% Shared layout adapted from the instructor's ME 537 homework sources.
\documentclass[12pt,letterpaper]{article}
\usepackage{mathptmx}
\usepackage[T1]{fontenc}
\usepackage{xcolor,amsmath,amssymb,xurl,listings,enumitem,tikz}
\textwidth=6.5in
\textheight=9in
\oddsidemargin=0in
\evensidemargin=0in
\topmargin=0in
\headheight=0in
\headsep=0in
\parskip=0.08in
\parindent=0in
\setlist[enumerate]{itemsep=0.04in,topsep=0.04in}
\lstset{basicstyle=\ttfamily\footnotesize,breaklines=true,columns=fullflexible,
keepspaces=true,showstringspaces=false,aboveskip=0.04in,belowskip=0.04in}
\newcommand{\R}[2]{{}^{#1}\!R_{#2}}
\newcommand{\T}[2]{{}^{#1}\!T_{#2}}
\newcommand{\courseheader}[2]{%
\setlength{\fboxrule}{2pt}\setlength{\fboxsep}{2mm}%
\fcolorbox{black}{white}{\parbox{\dimexpr\linewidth-2\fboxsep-2\fboxrule\relax}{\centering\textbf{Introduction to Robotics Homework \##1 -- Draft}}}\par
\textbf{#2}\par}
\newcommand{\conventions}{Use right-handed rotations, column vectors, radians,
and meters. $\R{A}{B}$ maps coordinates from frame $B$ into frame $A$;
$\T{A}{B}$ also includes the origin of $B$ expressed in $A$.}
\newcommand{\submission}{\textbf{Submit:} labeled hand derivations and sketches,
your changed Python source files, the requested numerical/visual evidence,
and an AI-use statement (including ``none'' if applicable). Show your reasoning
before using code to check hand calculations. Public checks are intentionally
limited. The instructor supplies the deadline and submission destination.}

\begin{document}
\courseheader{11}{Reward shaping for interception}
\textbf{Schedule:} week 12, after the two RL lectures. You do not implement
PPO; you choose what it optimizes and compare the result with a controller
you write yourself.

\textbf{Task.} A two-link planar arm ($0.4$ and $0.3$ m, base at the origin)
must bring its tool to an object sliding toward $+x$ at $0.05$--$0.12$ m/s and
move with it. Success: the tool is within $2.5$ cm of the object and within
$5$ cm/s of its velocity for three consecutive control periods. The arm runs
as in the project: $50$ Hz control, joint speed $1.5$ rad/s, joint
acceleration $4$ rad/s$^2$, and one period of delay. An episode ends on
success (\emph{terminated}) or after $4$ s (\emph{truncated}). The
environment is \path{src/intro_robotics/rl_interception.py}.

\begin{enumerate}
\item \textbf{Specify the problem.}
\begin{enumerate}
\item The observation holds $\cos q$, $\sin q$, $\dot q$, the previous action,
and the object's position and velocity relative to the tool; the action is
two numbers in $[-1,1]$ that scale the joint speed limit. Why does the
observation include the previous action? (Hint: when is it applied?) Why
relative position rather than the object's world position?
\item Explain the difference between terminated and truncated, and why PPO
should not treat running out of time as reaching a goal.
\item For one step with distance $0.20$ m, relative speed $0.10$ m/s,
action $[0.6,-0.8]$, and no success, compute the three rewards defined in
problem 2. Repeat for distance $0.02$ m, relative speed $0.01$ m/s, and
success.
\end{enumerate}

\item \textbf{Write three rewards and a conventional baseline.}
In \path{homework/hw11/rewards.py}, complete
\begin{itemize}
\item \texttt{sparse\_reward}: $1$ on success, else $0$;
\item \texttt{distance\_reward}: $-d$ every step, plus $10$ on success;
\item \texttt{shaped\_reward}: \texttt{distance\_reward}
$-\,0.5\,\|v_{\rm rel}\|-0.01\,\|a\|^2$.
\end{itemize}
In \path{homework/hw11/baseline.py}, complete \texttt{InterceptController}:
with $J$ the first two rows of your \texttt{jacob\_planar}, estimate the
object's velocity $v_o=v_{\rm rel}+J\dot q$, command the tool velocity
$v_o+k\,(p_{\rm rel}+2T\,v_o)$ with $k=4$ s$^{-1}$ and $T=0.02$ s, convert
it to joint rates with your HW06 \texttt{damped\_least\_squares}, and scale
by the speed limit. Why the $2T$ lead? Run
\begin{lstlisting}
python -m unittest discover -s tests/public -p "test_hw11_*.py" -v
python -m homework.hw11.evaluate --policy baseline
\end{lstlisting}
\end{enumerate}

\newpage
\begin{enumerate}[start=3]
\item \textbf{Train and compare.} For each reward, train one policy:
\begin{lstlisting}
python -m homework.hw11.train --reward sparse --out tmp/hw11/sparse.zip
python -m homework.hw11.evaluate --policy tmp/hw11/sparse.zip
\end{lstlisting}
and the same for \texttt{distance} and \texttt{shaped}. Each run (600k steps)
takes about 5--10 minutes of unattended CPU time and prints it; on a much
slower computer use \texttt{-{}-steps 300000} and say so. Then train your
best reward again with \texttt{-{}-seed 2}.
\begin{enumerate}
\item Make one table: policy, training minutes, success out of 100,
mean time to intercept, and one observed failure (final distance and
relative speed). Include your baseline.
\item Which rewards learned the task? Explain the sparse result in terms of
what the policy experiences early in training.
\item Compare \texttt{distance} with \texttt{shaped}. Which failures did the
relative-speed term change? Did the action penalty change anything you can
measure?
\item Compare your two seeds of the same reward. What does the difference
say about reporting a single training run?
\end{enumerate}

\item \textbf{When is RL a plausible choice?}
\begin{enumerate}
\item List what your baseline needed from a model (kinematics, the delay,
the object's motion) and what the learned policy needed instead (a
simulator, a reward, and training time). Compare the time each controller
takes to intercept.
\item The success condition and the reward are different things. Give one
way a policy could earn a high return under \texttt{distance\_reward}
without intercepting, and how you would detect it.
\item Name one change (a sensor, the arm, the object's motion, or the
real robot) that would make you prefer the learned policy, and one that
would make you prefer the baseline. In the project, where could a policy
like this replace part of your controller, and what would it observe?
\end{enumerate}
\end{enumerate}

\submission
\end{document}
